\documentclass[12pt,a4paper]{article}

% ---- Encoding & Font ----
\usepackage{iftex}
\ifPDFTeX
  \usepackage[utf8]{inputenc}
  \usepackage[T1]{fontenc}
\else
  \usepackage{fontspec}  % XeTeX/LuaTeX (tectonic): Unicode nativo (·, —, ¿, ª, §)
\fi
\usepackage[spanish]{babel}
\usepackage{csquotes}

% ---- Page layout ----
\usepackage[top=2.5cm, bottom=2.5cm, left=2.5cm, right=2.5cm]{geometry}
\usepackage{setspace}
\onehalfspacing

% ---- Math ----
\usepackage{amsmath, amssymb, amsthm}
\usepackage{mathtools}
\usepackage{physics}
\usepackage{esint}  % \oiint

% ---- Graphics ----
\usepackage{graphicx}
\usepackage{tikz}
\usepackage{float}
\usepackage{caption}

% ---- Tables ----
\usepackage{array}
\usepackage{booktabs}
\usepackage{tabularx}
\usepackage{colortbl}

% ---- Colours ----
\usepackage{xcolor}
\definecolor{notebg}{HTML}{FFF3CD}
\definecolor{noteborder}{HTML}{FFC107}
\definecolor{boxred}{HTML}{F8D7DA}
\definecolor{boxredborder}{HTML}{F5C6CB}

% ---- Boxes ----
\usepackage{mdframed}
\usepackage{tcolorbox}
\tcbuselibrary{most}

\newtcolorbox{keybox}[1][]{
  colback=blue!5,
  colframe=blue!70!black,
  arc=4pt,
  boxrule=1pt,
  fonttitle=\bfseries,
  title={#1}
}

\newtcolorbox{warnbox}[1][]{
  colback=boxred,
  colframe=boxredborder,
  coltitle=black!75,
  arc=4pt,
  boxrule=1pt,
  fonttitle=\bfseries,
  title={#1}
}

\newtcolorbox{notebox}[1][]{
  colback=notebg,
  colframe=noteborder,
  coltitle=black!75,
  arc=4pt,
  boxrule=1pt,
  fonttitle=\bfseries,
  title={#1}
}

% ---- Hyperlinks & TOC ----
\usepackage[hidelinks]{hyperref}
\usepackage{bookmark}

% ---- Enumerate ----
\usepackage{enumitem}

% ---- Miscellaneous ----
\usepackage{icomma}
\usepackage{siunitx}
\usepackage{textcomp}
\usepackage[bottom]{footmisc}

% ---- Diagramas (gráficas y esquemas) ----
\usepackage{pgfplots}
\pgfplotsset{compat=1.18}
\usetikzlibrary{babel}  % babel español activa `>`; esto evita romper las flechas `->`
% courses/ElecMag2024/Clases/assets/tikzstyles.tex
% ---------------------------------------------------------------------------
% Estilos compartidos para los diagramas del curso Electromagnetismo (2024).
% Fuente única del "look" visual (colores, grosores, escalas) para que todas
% las figuras (TikZ / pgfplots / CircuiTikz) sean consistentes entre clases.
%
% Es una COPIA de courses/Fisica3-2015/Clases/assets/tikzstyles.tex: mismo
% dominio, misma paleta. Copia y no symlink para que cada curso sea
% autocontenido y la paleta de Física III pueda evolucionar aparte.
%
% Uso: la clase debe cargar antes `tikz` y `pgfplots` (y `circuitikz` si la
%      clase tiene circuitos), y luego % courses/ElecMag2024/Clases/assets/tikzstyles.tex
% ---------------------------------------------------------------------------
% Estilos compartidos para los diagramas del curso Electromagnetismo (2024).
% Fuente única del "look" visual (colores, grosores, escalas) para que todas
% las figuras (TikZ / pgfplots / CircuiTikz) sean consistentes entre clases.
%
% Es una COPIA de courses/Fisica3-2015/Clases/assets/tikzstyles.tex: mismo
% dominio, misma paleta. Copia y no symlink para que cada curso sea
% autocontenido y la paleta de Física III pueda evolucionar aparte.
%
% Uso: la clase debe cargar antes `tikz` y `pgfplots` (y `circuitikz` si la
%      clase tiene circuitos), y luego % courses/ElecMag2024/Clases/assets/tikzstyles.tex
% ---------------------------------------------------------------------------
% Estilos compartidos para los diagramas del curso Electromagnetismo (2024).
% Fuente única del "look" visual (colores, grosores, escalas) para que todas
% las figuras (TikZ / pgfplots / CircuiTikz) sean consistentes entre clases.
%
% Es una COPIA de courses/Fisica3-2015/Clases/assets/tikzstyles.tex: mismo
% dominio, misma paleta. Copia y no symlink para que cada curso sea
% autocontenido y la paleta de Física III pueda evolucionar aparte.
%
% Uso: la clase debe cargar antes `tikz` y `pgfplots` (y `circuitikz` si la
%      clase tiene circuitos), y luego \input{../assets/tikzstyles.tex}
% (compilar desde el directorio de la clase para que la ruta relativa resuelva).
% ---------------------------------------------------------------------------

% ---- Paleta (alineada con las cajas del preámbulo: keybox/notebox/warnbox) ----
\definecolor{figblue}{HTML}{1F4E79}   % azul principal (líneas, keybox)
\definecolor{figred}{HTML}{C0392B}    % rojo de énfasis (warnbox)
\definecolor{figamber}{HTML}{B8860B}  % ámbar (notebox)
\definecolor{figgray}{HTML}{555555}   % auxiliares, ejes, guías

% ---- CircuiTikz: escala y grosor de los componentes ----
% Guarda \ifdefined: la electrostática (Clases 1-14) no carga `circuitikz`, y
% sin la guarda el \input revienta con `Undefined control sequence \ctikzset`.
% Cuando el curso llegue a circuitos, el bloque se activa solo.
\ifdefined\ctikzset
  \ctikzset{
    bipoles/thickness=1,
    resistors/scale=0.9,
    capacitors/scale=0.9,
  }
\fi

% ---- TikZ: estilos reutilizables ----
% Nota: las puntas se declaran como `-{Latex[...]}` y no como `->`. La forma
% con llaves NO contiene un `>` literal, así que es inmune al `>` activo que
% introduce babel español (ver CLAUDE.md §7.3.3).
\usetikzlibrary{arrows.meta,calc,angles,quotes}

\tikzset{
  figline/.style={figblue, line width=0.9pt},
  figaux/.style={figgray, dashed, line width=0.6pt},
  % vectores
  figvec/.style={figblue, line width=0.9pt, -{Latex[length=2.2mm,width=1.6mm]}},
  figvecres/.style={figred, line width=1.3pt, -{Latex[length=2.6mm,width=1.9mm]}},
  figvecaux/.style={figgray, line width=0.7pt, -{Latex[length=1.8mm,width=1.3mm]}},
  figaxis/.style={figgray, line width=0.6pt, -{Latex[length=2mm,width=1.4mm]}},
  % cargas puntuales
  figcarga/.style={circle, inner sep=0pt, minimum size=5.4pt, draw=none},
  figpos/.style={figcarga, fill=figred},
  figneg/.style={figcarga, fill=figblue},
  figneutra/.style={figcarga, fill=figgray},
  % etiquetas
  figlbl/.style={font=\small},
  figlblaux/.style={font=\footnotesize, figgray},
}

% ---- pgfplots: estilo base de las gráficas ----
\pgfplotsset{
  emfig/.style={
    width=11cm, height=6.5cm,
    axis lines=left,
    xlabel style={font=\small},
    ylabel style={font=\small},
    tick label style={font=\footnotesize},
    every axis plot/.append style={line width=1pt},
    grid=major,
    grid style={figgray!20},
    legend cell align=left,
    legend style={font=\footnotesize, draw=figgray!40},
  },
}

% (compilar desde el directorio de la clase para que la ruta relativa resuelva).
% ---------------------------------------------------------------------------

% ---- Paleta (alineada con las cajas del preámbulo: keybox/notebox/warnbox) ----
\definecolor{figblue}{HTML}{1F4E79}   % azul principal (líneas, keybox)
\definecolor{figred}{HTML}{C0392B}    % rojo de énfasis (warnbox)
\definecolor{figamber}{HTML}{B8860B}  % ámbar (notebox)
\definecolor{figgray}{HTML}{555555}   % auxiliares, ejes, guías

% ---- CircuiTikz: escala y grosor de los componentes ----
% Guarda \ifdefined: la electrostática (Clases 1-14) no carga `circuitikz`, y
% sin la guarda el \input revienta con `Undefined control sequence \ctikzset`.
% Cuando el curso llegue a circuitos, el bloque se activa solo.
\ifdefined\ctikzset
  \ctikzset{
    bipoles/thickness=1,
    resistors/scale=0.9,
    capacitors/scale=0.9,
  }
\fi

% ---- TikZ: estilos reutilizables ----
% Nota: las puntas se declaran como `-{Latex[...]}` y no como `->`. La forma
% con llaves NO contiene un `>` literal, así que es inmune al `>` activo que
% introduce babel español (ver CLAUDE.md §7.3.3).
\usetikzlibrary{arrows.meta,calc,angles,quotes}

\tikzset{
  figline/.style={figblue, line width=0.9pt},
  figaux/.style={figgray, dashed, line width=0.6pt},
  % vectores
  figvec/.style={figblue, line width=0.9pt, -{Latex[length=2.2mm,width=1.6mm]}},
  figvecres/.style={figred, line width=1.3pt, -{Latex[length=2.6mm,width=1.9mm]}},
  figvecaux/.style={figgray, line width=0.7pt, -{Latex[length=1.8mm,width=1.3mm]}},
  figaxis/.style={figgray, line width=0.6pt, -{Latex[length=2mm,width=1.4mm]}},
  % cargas puntuales
  figcarga/.style={circle, inner sep=0pt, minimum size=5.4pt, draw=none},
  figpos/.style={figcarga, fill=figred},
  figneg/.style={figcarga, fill=figblue},
  figneutra/.style={figcarga, fill=figgray},
  % etiquetas
  figlbl/.style={font=\small},
  figlblaux/.style={font=\footnotesize, figgray},
}

% ---- pgfplots: estilo base de las gráficas ----
\pgfplotsset{
  emfig/.style={
    width=11cm, height=6.5cm,
    axis lines=left,
    xlabel style={font=\small},
    ylabel style={font=\small},
    tick label style={font=\footnotesize},
    every axis plot/.append style={line width=1pt},
    grid=major,
    grid style={figgray!20},
    legend cell align=left,
    legend style={font=\footnotesize, draw=figgray!40},
  },
}

% (compilar desde el directorio de la clase para que la ruta relativa resuelva).
% ---------------------------------------------------------------------------

% ---- Paleta (alineada con las cajas del preámbulo: keybox/notebox/warnbox) ----
\definecolor{figblue}{HTML}{1F4E79}   % azul principal (líneas, keybox)
\definecolor{figred}{HTML}{C0392B}    % rojo de énfasis (warnbox)
\definecolor{figamber}{HTML}{B8860B}  % ámbar (notebox)
\definecolor{figgray}{HTML}{555555}   % auxiliares, ejes, guías

% ---- CircuiTikz: escala y grosor de los componentes ----
% Guarda \ifdefined: la electrostática (Clases 1-14) no carga `circuitikz`, y
% sin la guarda el \input revienta con `Undefined control sequence \ctikzset`.
% Cuando el curso llegue a circuitos, el bloque se activa solo.
\ifdefined\ctikzset
  \ctikzset{
    bipoles/thickness=1,
    resistors/scale=0.9,
    capacitors/scale=0.9,
  }
\fi

% ---- TikZ: estilos reutilizables ----
% Nota: las puntas se declaran como `-{Latex[...]}` y no como `->`. La forma
% con llaves NO contiene un `>` literal, así que es inmune al `>` activo que
% introduce babel español (ver CLAUDE.md §7.3.3).
\usetikzlibrary{arrows.meta,calc,angles,quotes}

\tikzset{
  figline/.style={figblue, line width=0.9pt},
  figaux/.style={figgray, dashed, line width=0.6pt},
  % vectores
  figvec/.style={figblue, line width=0.9pt, -{Latex[length=2.2mm,width=1.6mm]}},
  figvecres/.style={figred, line width=1.3pt, -{Latex[length=2.6mm,width=1.9mm]}},
  figvecaux/.style={figgray, line width=0.7pt, -{Latex[length=1.8mm,width=1.3mm]}},
  figaxis/.style={figgray, line width=0.6pt, -{Latex[length=2mm,width=1.4mm]}},
  % cargas puntuales
  figcarga/.style={circle, inner sep=0pt, minimum size=5.4pt, draw=none},
  figpos/.style={figcarga, fill=figred},
  figneg/.style={figcarga, fill=figblue},
  figneutra/.style={figcarga, fill=figgray},
  % etiquetas
  figlbl/.style={font=\small},
  figlblaux/.style={font=\footnotesize, figgray},
}

% ---- pgfplots: estilo base de las gráficas ----
\pgfplotsset{
  emfig/.style={
    width=11cm, height=6.5cm,
    axis lines=left,
    xlabel style={font=\small},
    ylabel style={font=\small},
    tick label style={font=\footnotesize},
    every axis plot/.append style={line width=1pt},
    grid=major,
    grid style={figgray!20},
    legend cell align=left,
    legend style={font=\footnotesize, draw=figgray!40},
  },
}


% ---- Title ----
\title{\textbf{Clase 18: Función delta de Dirac y ley de Ampère} \\[0.3em]
        \large{Cargas puntuales como distribuciones y equivalencia entre Biot--Savart y Ampère en régimen estacionario}}
\author{Transcripto y expandido --- Curso Electromagnetismo (OpenFING)}
\date{}

\begin{document}

\maketitle
\thispagestyle{empty}
\newpage

\tableofcontents
\newpage

\section{Paréntesis: la delta de Dirac}

La \textbf{función delta de Dirac} se podría haber introducido en el primer capítulo
---ahí está en el libro---; el docente la postergó porque hasta ahora no hacía
falta y es un tema más técnico. A diferencia del paréntesis del símbolo de
Levi-Civita, que fue «por amor al arte», \textbf{esto forma parte del programa} y de
acá en más se usa.

\subsection{Qué es una carga puntual}

Como siempre que se trabaja con distribuciones continuas, se ignora la
cuantización de la carga (válido mientras las cargas en juego sean grandes
frente a la del electrón).

«Carga puntual» no quiere decir realmente puntual: quiere decir que su tamaño
es mucho menor que las distancias relevantes y, más precisamente, \textbf{mucho menor
que la resolución del aparato de medida}. Una distancia no es «chica» en
abstracto, es chica respecto de algo.

\begin{itemize}
  \item Dos bolitas separadas dos metros, medidas con resolución de un centímetro:
  son puntuales si son mucho más chicas que un centímetro.
  \item Si se quiere medir con un milímetro o una micra de precisión, el criterio se
  vuelve más exigente.
  \item El aparato tiene que ser apropiado al problema: una resolución de diez metros
  para cargas separadas dos metros no mide nada.
\end{itemize}

\begin{notebox}[Idea básica]
Si se ignora la cuantización de la carga, no es posible
diferenciar físicamente una carga puntual de una distribución continua de
carga \textbf{muy localizada}. De hecho esa es la definición de carga puntual que
se usa: una distribución localizada de modo que su tamaño sea mucho menor que
la resolución del aparato de medida.
\end{notebox}

\subsection{Distribuciones como límites de funciones}

Los matemáticos probaron que esto se puede construir de manera consistente: se
puede trabajar con algo que se parece a una distribución continua de carga pero
describe una carga puntual. La herramienta es la \textbf{teoría de distribuciones}:
objetos que no son funciones sino \textbf{límites de funciones}. Una carga cada vez
más concentrada no tiene por límite una función.

La teoría rigurosa lleva un curso entero y no se hace acá; se dan las ideas
intuitivas que se usan en física y cómo se manipula la distribución más simple,
la delta de Dirac. Dirac ---físico, años treinta del siglo pasado--- la introdujo
\textbf{antes} de que existiera la formulación matemática general.

\section{Regularización de la delta}

\subsection{Las dos propiedades}

Se empieza en una dimensión. Sobre una recta se toma una función
$\delta_\varepsilon(x)$ centrada en el origen, de ancho del orden de
$\sqrt\varepsilon$ y altura del orden de $1/\sqrt\varepsilon$. Cuando
$\varepsilon\to0$ se concentra, \textbf{pero el área no cambia}: se achica el ancho
y se agranda la altura. Pensada como densidad lineal de carga, la carga total
es la misma, cada vez más concentrada.

El aparato de medida se modela con una función $f(x)$, su \textbf{sensibilidad}:
mide la carga alrededor de una posición $x_0$ con un ancho de precisión
$\sqrt\Delta$. Para ese aparato la distribución se comporta como carga puntual
si $\varepsilon\ll\Delta$: no puede saber si es puntual o no.

Una \textbf{regularización} de la delta es una función suave que cumple:

\begin{keybox}
\[
\boxed{\ \begin{aligned}
&\text{(1)}\quad \delta_\varepsilon(x)\xrightarrow[\varepsilon\to0]{}0
\quad\text{para todo } x\neq0;\\[4pt]
&\text{(2)}\quad \int_{-\infty}^{+\infty} dx\;\delta_\varepsilon(x)\,f(x)
\xrightarrow[\varepsilon\to0]{} f(0).
\end{aligned}\ }
\]
\end{keybox}

La segunda dice que, aunque la función tiende a cero en todos lados salvo en el
origen, se concentra de tal manera que la integral contra la sensibilidad del
aparato sólo «conoce» el valor de $f$ en un punto. Hay muchas regularizaciones
posibles.

\subsection{Ejemplo: la gaussiana}

\[
\delta_\varepsilon(x) = \frac{1}{\sqrt{2\pi\varepsilon}}\;
\exp\!\left(-\frac{x^2}{2\varepsilon}\right).
\]

\begin{figure}[H]
  \centering
  % Clase 18 — regularización gaussiana de la delta de Dirac,
%   delta_eps(x) = exp(-x^2 / 2 eps) / sqrt(2 pi eps),
% para eps = 1, 0.25 y 0.05: ancho ~ sqrt(eps), altura ~ 1/sqrt(eps), área 1.
% En ámbar punteado, la sensibilidad f(x) de un aparato de medida (curva
% ilustrativa, centrada en x_0 = 0.6): cuando eps es mucho menor que su ancho,
% la integral sólo "ve" f(0).
\begin{tikzpicture}
  \begin{axis}[
    emfig, width=12cm, height=6.6cm,
    xmin=-3, xmax=3, ymin=0, ymax=1.95,
    xtick={-3,-2,-1,0,1,2,3}, ytick={0,0.5,1,1.5},
    xlabel={$x$}, ylabel={$\delta_\varepsilon(x)$},
    legend style={at={(0.5,-0.38)}, anchor=north, legend columns=4, /tikz/every even column/.append style={column sep=1em}},
  ]
    \addplot[figblue!45, domain=-3:3, samples=121] {exp(-x^2/2)/sqrt(2*pi)};
    \addlegendentry{$\varepsilon = 1$}
    \addplot[figblue!75, domain=-3:3, samples=161] {exp(-x^2/0.5)/sqrt(2*pi*0.25)};
    \addlegendentry{$\varepsilon = 0{,}25$}
    \addplot[figblue, domain=-1.5:1.5, samples=241] {exp(-x^2/0.1)/sqrt(2*pi*0.05)};
    \addlegendentry{$\varepsilon = 0{,}05$}
    \addplot[figamber, dashed, domain=-3:3, samples=121] {0.75*exp(-(x-0.6)^2/2.2)};
    \addlegendentry{$f(x)$}
  \end{axis}
\end{tikzpicture}

  \caption{La regularización gaussiana para tres valores de $\varepsilon$: se angosta y crece, con área constante. En punteado, una sensibilidad $f(x)$ ilustrativa: cuando $\varepsilon$ es mucho menor que su ancho, la integral sólo registra $f(0)$.}
\end{figure}

\textbf{Propiedad 1.} Inmediata: para cualquier $x\neq0$ la exponencial tiende a
cero mucho más rápido de lo que crece el prefactor.

\textbf{Propiedad 2.} Todo ocurre a distancia $\sqrt\varepsilon$ del origen, así
que se hace un zoom con el cambio $u = x/\sqrt\varepsilon$,
$du = dx/\sqrt\varepsilon$:

\[
\int_{-\infty}^{+\infty} dx\,f(x)\,
\frac{e^{-x^2/2\varepsilon}}{\sqrt{2\pi\varepsilon}}
= \int_{-\infty}^{+\infty} du\;f\!\left(\sqrt\varepsilon\,u\right)
\frac{e^{-u^2/2}}{\sqrt{2\pi}}
\;\xrightarrow[\varepsilon\to0]{}\;
f(0)\int_{-\infty}^{+\infty}\frac{e^{-u^2/2}}{\sqrt{2\pi}}\,du .
\]

Queda $f(0)$ por una integral que ya no depende de $\varepsilon$. La convención
es normalizar la delta para que el área total sea 1; el factor
$\sqrt{2\pi\varepsilon}$ estaba puesto para eso.

\begin{notebox}[Qué se le pide a $f$]
En la teoría de distribuciones se suele pedir que
sea $C^\infty$ y de soporte compacto (nula fuera de un intervalo finito). El
docente no es preciso al respecto y aclara que el argumento seguiría valiendo
pidiendo menos.
\end{notebox}

\subsection{La integral gaussiana}

Falta ver que $\int e^{-u^2/2}\,du = \sqrt{2\pi}$. «Lo que importa no es el
qué, sino el cómo.» El truco: un número positivo es la raíz de su cuadrado, y
el cuadrado es una integral doble en el plano $(u,v)$, que se pasa a polares
(el jacobiano trae un $r$):

\[
\int_{-\infty}^{+\infty} e^{-u^2/2}\,du
= \left[\int\!\!\int du\,dv\;e^{-(u^2+v^2)/2}\right]^{1/2}
= \left[\int_0^{2\pi}\!d\theta\int_0^{\infty}\! r\,e^{-r^2/2}\,dr\right]^{1/2}.
\]

La integral angular da $2\pi$. La radial, con $z = r^2/2$, $dz = r\,dr$, es
$\int_0^\infty e^{-z}\,dz = 1$. Por lo tanto vale $\sqrt{2\pi}$ y

\[
\int_{-\infty}^{+\infty} dx\;\delta_\varepsilon(x)\,f(x)\to f(0).
\]

\subsection{El objeto límite}

\begin{notebox}[$\delta_\varepsilon(x)$ no tiene límite como función\textbf{ ---se vuelve
infinitamente alta e infinitamente angosta---, }pero el proceso de integrarla
contra funciones suaves sí tiene límite]
Esa operación límite es lo que se
llama una \textbf{distribución}.
\end{notebox}

Al tomar el límite se omite el $\varepsilon$: el objeto es la \textbf{delta de
Dirac}, $\delta(x)$, y corresponde a la carga puntual. Para el curso el rigor
del límite no importa; alcanza con que las dos descripciones sean físicamente
equivalentes.

\section{La delta en tres dimensiones y las cargas puntuales}

\subsection{Densidad de una carga puntual}

La densidad asociada a una carga puntual $q$ en el origen es

\[
\rho(\vec r) = q\,\delta(\vec r),
\qquad
\delta(\vec r) \equiv \delta(x)\,\delta(y)\,\delta(z) :
\]
el producto de las tres deltas, que localiza en $x=0$, $y=0$ y $z=0$ a la vez.
Para una carga en $\vec r_0$:

\begin{keybox}
\[
\boxed{\ \rho(\vec r) = q\,\delta(\vec r-\vec r_0)\ }
\]
\end{keybox}
concentrada donde se anula el argumento. Integrando contra una función suave,
con el cambio $\vec r\,' = \vec r - \vec r_0$ (que no altera el volumen):

\[
\int_V f(\vec r)\,q\,\delta(\vec r-\vec r_0)\,dV
= q\int_V f(\vec r_0+\vec r\,')\,\delta(\vec r\,')\,dV' = q\,f(\vec r_0).
\]

\begin{notebox}[Regla operativa]
Al integrar, la delta \textbf{elimina la integral y evalúa la
variable de integración en el punto donde se anula su argumento.}
\end{notebox}

De acá en más las cargas puntuales y las distribuciones continuas se tratan al
mismo nivel; no hace falta distinguir casos.

\subsection{Ejemplo: potencial de un conjunto de cargas}

Chequeo de coherencia. El potencial de una distribución continua es

\[
\phi(\vec r) = \frac{1}{4\pi\varepsilon_0}\int_V
\frac{\rho(\vec r\,')}{|\vec r-\vec r\,'|}\,dV' .
\]

Para cargas $q_i$ en posiciones $\vec r_i$ se pone
$\rho(\vec r\,') = \sum_i q_i\,\delta(\vec r\,'-\vec r_i)$; cada delta evalúa
el integrando en $\vec r\,' = \vec r_i$:

\[
\phi(\vec r) = \frac{1}{4\pi\varepsilon_0}\sum_i q_i\int_V
\frac{\delta(\vec r\,'-\vec r_i)}{|\vec r-\vec r\,'|}\,dV'
= \frac{1}{4\pi\varepsilon_0}\sum_i\frac{q_i}{|\vec r-\vec r_i|},
\]
que es el potencial de una suma de cargas puntuales, como tenía que ser.

\section{Una identidad clave: la divergencia de \texorpdfstring{$\vec r/r^3$}{r/r3}}

\subsection{De dónde sale}

Se combinan dos cosas que ya se saben. La ley de Gauss,
$\div \vec E = \rho/\varepsilon_0$, con $\rho = q\,\delta(\vec r)$ ---una
carga puntual en el origen---, y el campo de esa carga, que es Coulomb,
$\vec E = q\,\vec r/(4\pi\varepsilon_0 r^3)$. Si la delta tiene sentido, las dos
tienen que ser coherentes:

\[
\div \left(\frac{1}{4\pi\varepsilon_0}\frac{q\,\vec r}{r^3}\right)
= \frac{q}{\varepsilon_0}\,\delta(\vec r)
\quad\Longrightarrow\quad
\boxed{\ \div \left(\frac{\vec r}{r^3}\right) = 4\pi\,\delta(\vec r)\ }
\]

«Una fórmula maravillosa», que se va a usar mucho. Se verifica que hace lo que
una delta debe hacer.

\subsection{Verificación fuera del origen}

Tiene que valer cero para $\vec r\neq0$. En cartesianas, con
$r = (x^2+y^2+z^2)^{1/2}$:

\[
\pdv{x}\,\frac{x}{r^3} = \frac{1}{r^3} - \frac{3x^2}{r^5},
\]
y lo análogo para $y$ y $z$. Sumando los tres:

\[
\div \frac{\vec r}{r^3} = \frac{3}{r^3} - \frac{3\,(x^2+y^2+z^2)}{r^5} = 0
\qquad (\vec r\neq0).
\]

\begin{warnbox}
La cuenta \textbf{no vale en el origen}: ahí la función no es diferenciable.
\end{warnbox}

\subsection{Verificación de la propiedad integral}

Hay que ver que $\int \div (\vec r/r^3)\,f\,dV = 4\pi f(0)$, resultado
raro para algo que es cero en todos lados salvo en un punto. Se usa «el viejo
truco de sumar y restar», la derivada de un producto:

\[
\div \!\left(\frac{\vec r}{r^3}\right) f
= \div \!\left(\frac{\vec r}{r^3}\,f\right) - \frac{\vec r}{r^3}\cdot\vnabla f .
\]

Al primer término se le aplica el teorema de la divergencia; en el segundo,
$\vec r/r^3$ es radial y sólo sobrevive la componente radial del gradiente:

\[
\int_V \div \!\left(\frac{\vec r}{r^3}\right) f\,dV
= \oiint_S \frac{\vec r\cdot\hat n}{r^3}\,f\,dS
- \int_V \frac{1}{r^2}\,\pdv{f}{r}\,dV .
\]

Se toma como $V$ una esfera de radio $R$ centrada en el origen (si el origen no
estuviera adentro, daría cero). En esféricas el $r^2$ del elemento de volumen
cancela el $1/r^2$ y la integral radial es inmediata:

\[
\int_V \frac{1}{r^2}\,\pdv{f}{r}\,dV
= \int_0^{\pi}\!d\theta\,\sin\theta\int_0^{2\pi}\!d\varphi\,
\Big[f(R,\theta,\varphi) - f(0)\Big].
\]

Haciendo $R\to\infty$: $f$ es la sensibilidad de un aparato, está localizada,
y $f(R,\theta,\varphi)\to0$; por lo mismo se anula el término de superficie.
En $r=0$ ya no hay dependencia angular ---es un punto--- y la integral angular da
$4\pi$:

\[
\int \div \!\left(\frac{\vec r}{r^3}\right) f\,dV = 4\pi\,f(0).\ \blacksquare
\]

\begin{warnbox}[Hipótesis usadas]
$f$ suave (para poder derivar) y $f(\vec r)\to0$
suficientemente rápido cuando $r\to\infty$; y hubo que integrar sobre todo el
espacio.
\end{warnbox}

Cuando Dirac la introdujo, la reacción fue «nos está tomando el pelo». Después
los matemáticos lo emprolijaron: hay que pensar las distribuciones como la
operación de integrar contra una función en un límite; en la práctica siempre
hay una integral de por medio, y ahí sí hay límite.

\section{Ley de Ampère}

\subsection{Qué se va a probar y cuándo vale}

En Física III Biot--Savart y Ampère se usaban como equivalentes sin probarlo.
Probarlo sin la delta se puede, pero es mucho más trabajoso.

\begin{warnbox}[¿Equivalentes?]
Al final de Física III se ve que la ley de Ampère es
matemáticamente inconsistente para corrientes no estacionarias (el problema
que resuelve el término de Maxwell). Nada puede ser equivalente a algo
inconsistente; la equivalencia vale \textbf{para corrientes estacionarias}.
\end{warnbox}

Régimen estacionario: la densidad de carga no depende del tiempo (no es que sea
cero), y la ecuación de continuidad
$\partial\rho/\partial t + \div \vec J = 0$ se reduce a

\[
\div \vec J = 0 .
\]

Ya se vio que Biot--Savart implica $\div \vec B = 0$. Ahora se busca el
\textbf{rotacional}, respecto de la variable del punto de campo, $\vnabla_2\times$,
partiendo de

\[
\vec B(\vec r_2) = \frac{\mu_0}{4\pi}\int_V
\vec J(\vec r_1)\times\frac{\vec r_2-\vec r_1}{|\vec r_2-\vec r_1|^3}\,dV_1 .
\]

\subsection{Rotacional del integrando}

Fórmula de cálculo vectorial, que queda como \textbf{ejercicio} (con Levi-Civita o
«a manopla»):

\[
\curl (\vec F\times\vec G)
= \vec F\,(\div \vec G) - \vec G\,(\div \vec F)
+ (\vec G\cdot\vnabla)\vec F - (\vec F\cdot\vnabla)\vec G .
\]

Con $\vec F = \vec J(\vec r_1)$, que \textbf{no depende de $\vec r_2$}, y
$\vec G = (\vec r_2-\vec r_1)/|\vec r_2-\vec r_1|^3$, todas las derivadas que
caen sobre $\vec F$ son nulas. Sobreviven dos términos:

\[
\vnabla_2\times\left(\vec J(\vec r_1)\times\vec G\right)
= \vec J(\vec r_1)\;\vnabla_2\cdot\vec G - \left(\vec J(\vec r_1)\cdot\vnabla_2\right)\vec G .
\]

\begin{itemize}
  \item El primero es la identidad de §4:
  $\vnabla_2\cdot\vec G = 4\pi\,\delta(\vec r_2-\vec r_1)$.
  \item En el segundo, $\vec G$ depende sólo de $\vec r_2-\vec r_1$, así que derivar
  respecto de $\vec r_2$ o de $\vec r_1$ difiere en un signo:
  $-(\vec J\cdot\vnabla_2)\vec G = +(\vec J\cdot\vnabla_1)\vec G$.
\end{itemize}

Hasta acá es sólo cálculo vectorial; todavía no se usó el régimen estacionario.

\[
\vnabla_2\times\vec B(\vec r_2) = \frac{\mu_0}{4\pi}\int_V\left[
4\pi\,\delta(\vec r_2-\vec r_1)\,\vec J(\vec r_1)
+ \left(\vec J(\vec r_1)\cdot\vnabla_1\right)
\frac{\vec r_2-\vec r_1}{|\vec r_2-\vec r_1|^3}\right]dV_1 .
\]

\subsection{El término con la delta}

Es el fácil: la delta elimina la integral y evalúa $\vec J$ en
$\vec r_1 = \vec r_2$; los $4\pi$ se cancelan:

\[
\frac{\mu_0}{4\pi}\int_V 4\pi\,\delta(\vec r_2-\vec r_1)\,\vec J(\vec r_1)\,dV_1
= \mu_0\,\vec J(\vec r_2).
\]

\subsection{El término de borde}

Se hace componente a componente; para la componente $x$, con
$G_x = (x_2-x_1)/|\vec r_2-\vec r_1|^3$. La idea es convertirlo en una
divergencia:

\[
\vnabla_1\cdot\left(\vec J(\vec r_1)\,G_x\right)
= \underbrace{\left(\vnabla_1\cdot\vec J(\vec r_1)\right)}_{=\,0}G_x
+ \vec J(\vec r_1)\cdot\vnabla_1 G_x .
\]

\textbf{Acá entra la física:} $\vnabla_1\cdot\vec J = 0$ por régimen estacionario.
Entonces el integrando es una divergencia pura y, por el teorema de la
divergencia,

\[
\int_V \vec J(\vec r_1)\cdot\vnabla_1 G_x\,dV_1
= \oiint_S G_x\;\vec J(\vec r_1)\cdot\hat n\,dS = 0,
\]
porque si las corrientes están \textbf{localizadas} basta tomar un volumen que las
contenga enteramente: sobre su borde $\vec J = 0$ (ni siquiera hace falta ir al
infinito). Lo mismo para $y$ y $z$.

\begin{keybox}
\[
\boxed{\ \curl \vec B = \mu_0\,\vec J\ }
\qquad\text{(corrientes localizadas y estacionarias)}
\]
\end{keybox}

Es la \textbf{ley de Ampère} en forma diferencial. Lo que se probó es que Biot y
Savart \textbf{implica} Ampère bajo esas dos hipótesis.

\subsection{Forma integral}

\begin{figure}[H]
  \centering
  % Clase 18 — forma integral de la ley de Ampère. Curva cerrada C que delimita
% una superficie S, orientada con la regla de la mano derecha. La circulación
% de B en C es mu_0 por la corriente que cruza S.
\begin{tikzpicture}[scale=1]

  % superficie (casquete)
  \fill[figblue!12] (-2.6,0) arc (180:360:2.6 and 0.85) .. controls (2.4,2.3) and (-2.4,2.3) .. (-2.6,0);
  \draw[figline] (2.6,0) .. controls (2.4,2.3) and (-2.4,2.3) .. (-2.6,0);

  % curva C: mitad de atras punteada, mitad de adelante llena
  \draw[figaux] (2.6,0) arc (0:180:2.6 and 0.85);
  \draw[figblue, line width=1.3pt] (-2.6,0) arc (180:360:2.6 and 0.85);
  \draw[figvecres] (-0.3,-0.85) -- (0.4,-0.85);
  \node[figlbl] at (2.95,-0.45) {$C$};
  \node[figlbl] at (-1.95,1.45) {$S$};

  % normales
  \draw[figvecaux] (-1.35,1.25) -- (-1.75,1.95);
  \draw[figvecaux] (1.35,1.25) -- (1.75,1.95);
  \node[figlblaux] at (2.02,2.0) {$\hat n$};

  % corriente que cruza S
  \foreach \jx in {-0.6,0,0.6} {
    \draw[figvecres] (\jx,-0.35) -- (\jx,2.75);
  }
  \node[figlbl, text=figred] at (0.95,2.75) {$\vec J$};

\end{tikzpicture}

  \caption{Curva $C$ y superficie $S$ que se apoya en ella, orientadas con la regla de la mano derecha. La circulación de $\vec B$ en $C$ es $\mu_0$ por la corriente que cruza $S$.}
\end{figure}

Para recuperar la de Física III se toma una curva cerrada $C$ que delimita una
superficie $S$, orientada con la regla de la mano derecha, y se integra el
flujo de ambos lados. Por el teorema de Stokes:

\[
\int_S (\curl \vec B)\cdot\hat n\,dS = \oint_C \vec B\cdot d\vec l,
\qquad
\mu_0\int_S \vec J\cdot\hat n\,dS = \mu_0\,I_S,
\]

\begin{keybox}
\[
\boxed{\ \oint_C \vec B\cdot d\vec l = \mu_0\,I_S\ }
\]
\end{keybox}
con $I_S$ la corriente que cruza $S$.

\section{Ejemplo: cable infinito}

Cable recto infinito con corriente estacionaria $I$.

\begin{warnbox}[Advertencia]
Esta corriente \textbf{no está localizada}, así que en rigor lo
probado no se aplica. Hay que pensarlo como caso límite de una corriente
localizada; se sabe que el límite está bien porque el resultado ya se obtuvo
con Biot y Savart.
\end{warnbox}

\begin{figure}[H]
  \centering
  % Clase 18 — cable infinito y curva amperiana circular en un plano
% perpendicular. Las líneas de B son cerradas (no hay monopolos) y el problema
% es invariante por rotaciones: son círculos, y |B| es constante sobre cada uno.
\begin{tikzpicture}[scale=1]

  % cable, tramo de abajo
  \draw[figblue, line width=1.8pt] (0,-2.1) -- (0,-0.62);

  % plano
  \fill[figgray!12] (-3.4,-0.9) -- (2.2,-0.9) -- (3.4,0.9) -- (-2.2,0.9) -- cycle;
  \draw[figgray, line width=0.6pt] (-3.4,-0.9) -- (2.2,-0.9) -- (3.4,0.9) -- (-2.2,0.9) -- cycle;

  % curva amperiana
  \draw[figred, line width=1pt, dashed] (0,0) ellipse (1.9 and 0.6);
  \draw[figaux] (0,0) -- (1.9,0);
  \node[figlblaux] at (0.95,0.17) {$r$};

  % campo tangente (sentido antihorario visto desde arriba)
  \draw[figvec] (1.9,0) -- ++(0,0.62);
  \draw[figvec] (-1.9,0) -- ++(0,-0.62);
  \draw[figvec] (0,0.6) -- ++(-0.95,0);
  \draw[figvec] (0,-0.6) -- ++(0.95,0);
  \node[figlbl] at (2.25,0.55) {$\vec B$};

  % cable, tramo de arriba
  \draw[figblue, line width=1.8pt, -{Latex[length=2.8mm,width=2mm]}] (0,0) -- (0,2.3);
  \fill[figblue] (0,0) circle (1.6pt);
  \node[figlbl] at (0.3,2.0) {$I$};

\end{tikzpicture}

  \caption{Cable infinito y curva amperiana circular en un plano perpendicular. El campo es tangente al círculo y de módulo constante sobre él.}
\end{figure}

Curva amperiana: un círculo de radio $r$ en un plano perpendicular al cable.
¿Por qué el campo es tangencial? Las líneas de $\vec B$ tienen que ser cerradas
(no hay monopolos) y el problema tiene simetría de rotación: las únicas curvas
cerradas con esa simetría son círculos. Además el módulo es constante sobre el
círculo:

\[
B\,2\pi r = \mu_0 I
\quad\Longrightarrow\quad
B = \frac{\mu_0 I}{2\pi r},
\]
el mismo resultado de la clase anterior. Lo que con Biot y Savart llevó un
pizarrón entero, con Ampère lleva una línea.

\begin{center}
\renewcommand{\arraystretch}{1.35}
\small
\begin{tabularx}{\textwidth}{l>{\raggedright\arraybackslash}X}
\toprule
 & \textbf{Cuándo conviene} \\
\midrule
Ampère & Problemas con mucha simetría \\
Biot y Savart & Poca simetría: hay que hacer las integrales, analítica o numéricamente; siempre funciona \\
\bottomrule
\end{tabularx}
\end{center}

\noindent\emph{La próxima clase: algún ejemplo más de Ampère y el \textbf{potencial vector}. Como
$\curl \vec B\neq0$ donde hay corriente, no existe en general un potencial
escalar magnético; pero sí algo parecido a un potencial, que es un vector.}

\end{document}
